% ============================================================
% templates/t00-core.tex   —— 底座：骨架 + T1 纯文字分组
%
% 这个文件是"所有版式的公共底座"。原语本身定义在 primitives.tex 里，
% 本文件只登记 T1 的适用场景与容量，并说明为什么原语不放在这里。
%
% 【T1 纯文字分组】
% 适用场景：一段说明、几条并列的事实、任何"就是几行字"的页。
%          没有标签、没有表格、没有图的页，默认都是它。
% 提供    ：\DeckLine \DeckGap（定义在 primitives.tex）
% 参数    ：无。每行一个 \DeckLine
% 容量    ：一页最多 \SlideMaxRows 行（推荐值，定义在 styles/ 里；
%           几何硬上限是 8 行，见 styles/default.tex 末尾的实测段）
%           一行最多 \SlideLineUnitsSafe 个全角字
%           （西文数字约占半角；几何硬上限 \SlideLineUnits 个）
% 分组    ：需要分段就插一个 \DeckGap。这是唯一的分组手段 ——
%           不要用分隔线、小标题或色块。
% 层级    ：只有"加粗"和"留白"两种手段，整页字号完全一致。
% 易错    ：折行点不受控。写文案前先按 \SlideLineUnitsSafe 数字数，
%           超了就压短，别指望它自己排好看。
%
% 【填满的样子】
%   \DeckLine{每一条要点写成一行，行与行之间自动留出间距。}
%   \DeckLine{想强调某个词时用 \textbf{加粗}，一句话里不要超过一处。}
%   \DeckLine{整页字号、字重、颜色完全一致，层级靠留白分组。}
%
%   \DeckGap
%
%   \DeckLine{这是第二组内容。上面那个 \DeckGap 就是分组手段。}
%
% 【为什么原语不定义在这里】\DeckLine 这类原语被 T2、T5 等模板调用。
% 如果它定义在 T1 里，那么"所有模板都依赖 T1"—— 删掉 T1 就会连带
% 弄坏别的模板。放在 primitives.tex 里则依赖是单向的：模板 → 工具箱。
% ============================================================

% T1 不需要新命令：它就是 \DeckLine / \DeckGap 的直接使用。
